% =====================================================================
%  Chapter 01 — Origin: the T-first thesis
% =====================================================================
\chapter{Origin: Temperature as the Master Variable}
\label{ch:origin}

\section{The problem}
\label{sec:the-problem}

The Clay Millennium Prize problem for the three-dimensional incompressible
Navier--Stokes equations asks whether smooth, divergence-free initial data of
finite energy generates a solution that stays smooth for all time. On the
torus \(\T^3=\R^3/(2\pi\Z)^3\) --- the setting this programme works in
throughout --- the equations are
\begin{equation}\label{eq:NS-prize}
  \partial_t u + (u\cdot\nabla)u = -\nabla p + \nu\Delta u,
  \qquad \nabla\cdot u = 0,
  \qquad u(\cdot,0)=u_0 .
\end{equation}
Leray's 1934 construction gives a global weak solution; Prodi and Serrin give
criteria under which a weak solution is smooth; Caffarelli, Kohn and Nirenberg
bound the parabolic Hausdorff dimension of the possible singular set. What is
missing, and has been missing for ninety years, is a mechanism that forbids
concentration at the critical scaling.

The obstruction is scale-criticality, and it is easy to state. Under the
scaling \(u_\lambda(x,t)=\lambda u(\lambda x,\lambda^2t)\), which preserves
\eqref{eq:NS-prize}, the energy \(\norm{u}_{L^2}^2\) is \emph{supercritical}:
it grows as \(\lambda^{-1}\) and therefore says nothing at small scales. Every
a priori quantity the equations hand over for free is either supercritical or
exactly critical, and the Prodi--Serrin, Ladyzhenskaya--Prodi--Serrin (LPS)
family of sufficient conditions sits precisely on the critical line, one
endpoint out of reach.

\section{The thesis}
\label{sec:the-thesis}

The T-First programme began from a physical reading of that gap. Its
founding claim, stated in the programme's specification document
\texttt{tfirst\_prd\_v4.docx} and restated at the head of every working
session, is this:

\begin{record}
\textbf{The T-first thesis.} Temperature \(T(x,t)\) is a \emph{scalar} field,
and it is the master variable. Every other fluid quantity --- density
\(\rho\), pressure \(p\), dynamic viscosity \(\mu\), thermal conductivity
\(k\), and velocity \(u\) itself --- is downstream of \(T\). The LPS
regularity gap is a thermodynamic artefact, and it closes through
\[
  A(T) \;=\; \frac{k(T)}{\rho(T)\,c_v(T)} \;>\; 0 ,
\]
a direct consequence of the second law of thermodynamics. Because \(T\)
satisfies a quasilinear \emph{scalar} parabolic equation, the strong maximum
principle applies to it --- a tool that is simply unavailable to any approach
working directly with the vector field \(u\) or the vorticity \(\omega\).
\end{record}

The three ingredients deserve separating, because they are of very different
kinds and they fared very differently.

\subsection{Scalarity}

The first ingredient is structural. A quasilinear parabolic equation for a
scalar unknown admits the strong maximum principle, the Harnack inequality,
De~Giorgi--Nash--Moser regularity, and a comparison principle. None of these
has a vector-valued analogue of comparable strength: for systems, the maximum
principle fails in general, and this failure is one of the reasons
three-dimensional Navier--Stokes is hard while its two-dimensional counterpart
is not. If the essential dynamics could be routed through a scalar, the
scalar's toolbox would become available.

This is a real and correct observation, and it is worth recording that it
survives everything that follows. What did not survive is the specific
scalar. The programme's later analytical work (Part~IV) does route the
argument through scalar quantities --- the vorticity magnitude
\(\abs{\omega}\), the coherence deficit \(\sstar\), the auxiliary scalar
\(\theta\) --- and does draw exactly the advantages the thesis predicted:
\texttt{lem:direction-eq} (line~5351) is derived with no pressure in it, which
is the structural benefit the vorticity route had been seeking. But by then
the master scalar was not \(T\).

\subsection{Downstream ordering}

The second ingredient is a claim about causation in the constitutive
relations, and it is illustrated in Figure~\ref{fig:downstream}. The
programme's own architecture rules encode it: \(T\) is the primary variable;
the strain rate \(S_{ij}\), the viscous stress \(\tau_{ij}=2\mu(T)S_{ij}\) and
the heat flux \(q\) are derived tensors and vectors, and are forbidden from
appearing as primary variables in the master equation.

\begin{figure}[ht]
\centering
\begin{tikzpicture}[
  >={Latex[length=2mm]},
  prim/.style={draw, thick, circle, minimum size=13mm, fill=panelgrey,
               font=\small\bfseries},
  der/.style={draw, rounded corners=2pt, inner sep=3.5pt, font=\small,
              minimum width=17mm, align=center},
  ten/.style={draw, dashed, rounded corners=2pt, inner sep=3.5pt,
              font=\small, minimum width=17mm, align=center},
  ar/.style={->, shorten >=1pt}
]
\node[prim] (T) at (0,0) {\(T\)};

\node[der] (mu)  at (3.2, 2.2) {\(\mu(T)\)};
\node[der] (k)   at (3.2, 1.1) {\(k(T)\)};
\node[der] (A)   at (3.2, 0.0) {\(A(T)=\dfrac{k}{\rho c_v}\)};
\node[der] (rho) at (3.2,-1.3) {\(\rho,\;c_v\)};
\node[der] (p)   at (3.2,-2.3) {\(p\)};

\node[der] (u)   at (7.2, 0.9) {\(u\)};
\node[ten] (S)   at (7.2,-0.4) {\(S_{ij}\)};
\node[ten] (tau) at (7.2,-1.7) {\(\tau_{ij}=2\mu S_{ij}\)};

\draw[ar] (T) -- (mu);
\draw[ar] (T) -- (k);
\draw[ar] (T) -- (A);
\draw[ar] (T) -- (rho);
\draw[ar] (T) -- (p);
\draw[ar] (mu) -- (u);
\draw[ar] (A)  -- (u);
\draw[ar] (u)  -- (S);
\draw[ar] (S)  -- (tau);
\draw[ar] (mu) to[bend left=18] (tau);

\begin{scope}[on background layer]
  \node[draw=rulegrey, dashed, rounded corners=3pt, fit=(S)(tau),
        inner sep=5pt, label={[font=\scriptsize\itshape,
        text=rulegrey]below:downstream: never primary}] {};
\end{scope}

\node[font=\scriptsize\itshape, text=rulegrey, align=center]
  at (0,-2.0) {master\\scalar};
\end{tikzpicture}
\caption[The downstream ordering of the T-first thesis]{The downstream
ordering asserted by the T-first thesis. \(T\) is primary; the transport
coefficients and the equation of state are functions of \(T\); the velocity is
driven by them; and the strain rate and viscous stress are derived from the
velocity. The programme's architecture rules forbid \(S_{ij}\) and
\(\tau_{ij}\) from ever entering the master equation as primary unknowns.}
\label{fig:downstream}
\end{figure}

\subsection{Positivity from the second law}

The third ingredient is the one that carried the physical weight. Thermodynamic
stability requires \(k>0\) and \(c_v>0\); with \(\rho>0\) this gives
\(A(T)>0\) pointwise, unconditionally, for any admissible fluid. The
programme treated this as an axiom to be \emph{verified numerically at every
step} rather than assumed: the property engine's \(A\)-field routine asserts
strict positivity on every call in debug mode, and a second-law check must
pass for a fluid before any solver run.

The payoff the thesis promised is a continuation criterion. If \(T\) is
governed by a quasilinear parabolic equation with a uniformly positive
diffusivity, then \(T\) obeys a maximum principle, stays in a compact interval
\([T_{\min},T_{\max}]\), and hence \(\mu(T)\) and \(A(T)\) are uniformly
bounded above and below. In the programme's Phase~0 formulation: \emph{a
singularity requires \(A(T)\to0\) or \(\mu(T)\to0\)}. If neither can happen,
neither can the singularity.

\section{Route~1: the system the thesis generates}
\label{sec:route1-system}

Made concrete in Prize geometry --- constant density, exactly
divergence-free velocity --- the thesis produces the following system,
called Route~1 throughout the programme:
\begin{align}
  \rho\bigl(\partial_t u + (u\cdot\nabla)u\bigr)
    &= -\nabla p + \nabla\cdot\bigl(\mu(T)\nabla u\bigr),
    \qquad \nabla\cdot u = 0, \label{eq:route1-u}\\[2pt]
  \rho\,c_v\,\partial_t T
    &= \nabla\cdot\bigl(k(T)\nabla T\bigr)
       - \rho\,c_v\,(u\cdot\nabla)T
       + \mu(T)\abs{\nabla u}^2, \label{eq:route1-T}
\end{align}
with the ideal-gas-like closures used throughout the numerical work,
\[
  \mu(T)=\mu_0\Bigl(\frac{T}{T_0}\Bigr)^{3/2},\qquad
  k(T)=k_0\Bigl(\frac{T}{T_0}\Bigr)^{3/2},\qquad
  A(T)=\frac{k(T)}{\rho c_v}>0 .
\]
Equation~\eqref{eq:route1-T} is the master equation. It is scalar; its
diffusion coefficient is \(A(T)>0\); its source \(\mu(T)\abs{\nabla u}^2\) is
non-negative, being viscous dissipation; and the maximum principle therefore
gives a one-sided bound for free. The Prize equations \eqref{eq:NS-prize} are
recovered by \(\mu(T)\to\nu\), which is what happens when the temperature
field becomes uniform.

\section{What the thesis promised}
\label{sec:what-promised}

The programme laid out what it expected to obtain, in a form specific enough
to be falsified. Table~\ref{tab:promises} records the promises and, for the
reader's orientation, what actually became of each. The chapters that follow
justify the right-hand column; it is placed here so that the founding claims
are never read without their outcome attached.

\begin{table}[ht]
\centering
\small
\begin{tabular}{@{}p{0.44\textwidth} p{0.48\textwidth}@{}}
\toprule
Promised & Outcome \\
\midrule
\(T\) is a scalar, so the strong maximum principle applies where vector
methods have nothing.
& Correct as a structural observation, and it did eventually pay --- but
  through \(\abs{\omega}\), \(\sstar\) and \(\theta\), not through \(T\)
  (Chapters~\ref{ch:claim-a}, \ref{ch:inheritance}). \\[3pt]
\(A(T)>0\) from the second law, uniformly.
& Verified in every numerical run the programme performed;
  \(A_{\min}>0\) in all \(14\) Layer-2 experiments, all \(10\) Layer-3
  \(\mu\)-sweep runs, and all three adversarial \(3\)D runs at
  \(N=64,128,256\) (Chapters~\ref{ch:layer12}, \ref{ch:layer3}). \\[3pt]
Continuation criterion: a singularity forces \(A(T)\to0\) or
\(\mu(T)\to0\).
& Implemented as \texttt{layer4/continuation\_criterion.py} and never
  triggered in any run. It is a criterion about Route~1, not about
  \eqref{eq:NS-prize}, and no Route-1 global existence theorem was
  proved. \OPEN. \\[3pt]
Prize equations recovered by a simple limit \(\mu(T)\to\nu\), with the
regularity margin surviving.
& The limit is numerically \emph{regular}: the deviation scales linearly,
  \(\abs{\varepsilon(\delta)}\propto\delta^{0.98}\) in 3D, with no collapse
  (Chapter~\ref{ch:layer3}). No analytical counterpart was proved. \\[3pt]
Six phases to a complete proof over forty-eight months
(Chapter~\ref{ch:architecture}).
& Phases~0--5 were never executed as planned. The programme's actual path
  left Route~1 within eighteen sessions and worked thereafter inside the
  exact Prize equations. \\
\bottomrule
\end{tabular}
\caption[What the founding thesis promised, and what became of it]{The
founding promises and their outcomes. The programme's most durable results
came from a route the founding thesis treated as a backup.}
\label{tab:promises}
\end{table}

\section{An honest note on where the thesis went}
\label{sec:where-thesis-went}

It should be said plainly, at the outset, because it colours everything after
Chapter~\ref{ch:layer3}: the analytical line that produced this programme's
strongest results does not use temperature.

From \Sn{18} onwards the proof document works with the exact incompressible
Navier--Stokes equations \eqref{eq:NS-prize} at constant viscosity, and with
the auxiliary scalar
\[
  \partial_t\theta + u\cdot\nabla\theta
    = \nu\Delta\theta + \nu\abs{\nabla u}^2, \qquad \theta(\cdot,0)=0,
\]
which is Route~2. The word ``temperature'' does not appear in the proof
document's mathematics. Whether \(\theta\) \emph{is} temperature in disguise
is a question of interpretation --- it satisfies the transport--diffusion
equation of a passive scalar heated by viscous dissipation, which is exactly
what the internal-energy equation \eqref{eq:route1-T} reduces to at constant
coefficients --- but nothing in the later mathematics depends on reading it
that way, and this book does not lean on the reading.

What survived from the founding thesis, then, is a method and not a
mechanism: \emph{find a scalar quantity that carries the dynamics, and use the
tools that only scalars admit.} The programme applied that method four times
--- to \(T\) (Route~1), to \(\theta\) (Route~2), to the vorticity magnitude
\(\abs{\omega}\) in the magnitude Caccioppoli of \Sn{34}, and to the coherence
deficit \(\sstar\) throughout \Sn{32}--\Sn{39}. Only the last two are still
load-bearing. The thesis was, in the end, right about the shape of the tool
and wrong about which scalar to put in it.
